\documentclass[11pt,a4paper]{article}
\usepackage{brainteaser}

\begin{document}
\brainteasertitle{9}{Les deux urnes}

\begin{statementbox}
On dispose de $50$ boules rouges et $50$ boules bleues. Il faut répartir les
$100$ boules entre deux urnes, avec au moins une boule dans chaque urne.

On choisit ensuite une urne uniformément au hasard, puis une boule uniformément
dans cette urne. Comment répartir les boules pour maximiser la probabilité de
tirer une boule rouge? Quelle est la probabilité optimale?
\end{statementbox}

\begin{hintsbox}
\begin{enumerate}
  \item Tester une stratégie extrême: isoler une seule boule rouge.
  \item Si l'urne 1 contient $r$ rouges et $b$ bleues, écrire la probabilité
  de succès en fonction de $r$ et $b$.
  \item Pour $r$ fixé, étudier la dépendance en $b$. Pourquoi le maximum est-il
  atteint à une extrémité?
  \item Après avoir placé toutes les boules bleues dans une même urne, combien
  de rouges faut-il placer seules dans l'autre?
\end{enumerate}
\end{hintsbox}

\correction

Si l'urne 1 contient $r$ boules rouges et $b$ boules bleues, la probabilité
de tirer une rouge vaut
\[
P(r,b)=\frac12\left(
\frac r{r+b}+\frac{50-r}{100-r-b}
\right).
\]

Pour $r$ fixé, cette expression est convexe en $b$ sur son domaine. Son
maximum est donc atteint à une extrémité: il est optimal de ne pas partager
les boules bleues entre les deux urnes. À permutation près, mettons toutes les
bleues dans l'urne 2 et seulement $p$ rouges dans l'urne 1. Alors
\[
P(p)=\frac12\left(1+\frac{50-p}{100-p}\right).
\]
La fraction $(50-p)/(100-p)$ décroît avec $p$. Il faut donc prendre la plus
petite valeur admissible, $p=1$.

\begin{keyideabox}
La répartition optimale est
\[
\text{urne 1: une rouge},
\qquad
\text{urne 2: 49 rouges et 50 bleues}.
\]
\end{keyideabox}

La probabilité obtenue est
\[
\frac12\cdot1+\frac12\cdot\frac{49}{99}
=\boxed{\frac{74}{99}\simeq74{,}75\%}.
\]

\begin{resultbox}
Isoler une seule boule rouge maximise la probabilité de succès, qui vaut
$74/99$.
\end{resultbox}

\end{document}
